\section{Conclusion}
\label{sec:conclusion}

One pipeline, run unchanged, builds surrogates of a two-parameter black-hole map and a three-parameter neutron-star map. At the 95th percentile of the relative error, the black-hole mass reaches \resultsBhMassTest{} significant figures against a data ceiling of \resultsBhMassCeiling{}, and its charge \resultsBhChargeTest{} against \resultsBhChargeCeiling{}; the neutron-star mass reaches \resultsNsMassTest{} against \resultsNsMassCeiling{}, and its charge \resultsNsChargeTest{} against \resultsNsChargeCeiling{} (Section~\ref{sec:results}). No pair reaches its ceiling. The surrogates stay inside the physical domain, where stable solutions exist, and refuse an input outside it (Section~\ref{sec:validity}).

The surrogate reads the maximum mass of a neutron star at the stable edge of each curve; its 95th percentile error is \discussionMassPNinetyFiveDigits{} figures, short of the ceiling of \discussionMassCeiling{} (Section~\ref{sec:discussion}). How the data-guided hypotheses $H_1$, $H_2$ and $H_3$ held is in the appendix (Section~\ref{sec:hypotheses}).

Two extensions are left for future work: Markov chain Monte Carlo estimation of the parameters $(\beta, \lambda)$ from synthetic observations~\citep{Reed2024}, and rapidly rotating compact objects~\citep{Komatsu1989,Cook1994,Stergioulas1995}. The code is open source (Section~\ref{sec:availability}).

\section*{Code and Data Availability}
\label{sec:availability}

The code behind this paper -- the data preparation, the exploratory analysis, the surrogate pipeline and the paper itself -- is available at \paperrepo{} under the MIT licence; a single \texttt{make} rebuilds every figure and number from the datasets. The trained networks are not distributed as fixed artifacts: they are rebuilt from the data by the same pipeline. The black-hole and neutron-star datasets~\citep{Doneva2017_BH,Doneva2017_NS} are not public.\review{State how a reader obtains the datasets, e.g.\ from the authors on reasonable request.}
